\documentclass[11pt]{article}
\usepackage[T1]{fontenc}
\usepackage{lmodern,amsmath,amssymb,amsthm,booktabs,microtype,xcolor}
\usepackage[letterpaper,margin=1in]{geometry}
\usepackage[colorlinks=true,linkcolor=blue!45!black,urlcolor=blue!45!black]{hyperref}
\newtheorem{theorem}{Theorem}[section]
\newtheorem{lemma}[theorem]{Lemma}
\newtheorem{proposition}[theorem]{Proposition}
\newtheorem{corollary}[theorem]{Corollary}
\newtheorem{remark}[theorem]{Remark}
\newcommand{\PP}{\mathbb P}
\newcommand{\EE}{\mathbb E}
\newcommand{\Z}{\mathcal Z}
\newcommand{\cO}{\mathcal O}
\newcommand{\comp}{\operatorname{comp}}
\newcommand{\tr}{\operatorname{tr}}
\newcommand{\half}{\tfrac12}
\setlength{\emergencystretch}{3em}
\title{The unrestricted sharp-order cycle bound}
\author{Collaborative research working draft}
\date{11 November 2026}
\begin{document}
\maketitle
\begin{abstract}
Let $\psi(r)$ be the maximum number of cycles in a finite graph of cycle rank $r$.
We prove that an absolute constant $K$ satisfies
\[
 \psi(r)\le 2^{r-1}+Kr2^{r/2}\qquad(r\ge1).
\]
The known cyclic-square examples have excess $\Theta(r2^{r/2})$, so this error order is necessary.

The new step removes the dependence on a positive forest excess from repeated short-cycle packing.
Two disjoint clean regions have uniformly bounded conditional affine-forest probabilities.
Conditioning on the other edges makes these probabilities multiply.
This controls every stage of cycle removal and turns a nested-chain argument into an exact half bound at sufficiently large fixed girth.
The established sharp bound for cuts of size at most four, the bounded-girth deficit, and induction on rank complete the proof.

The local estimates use only the original qualitative marked-cycle forest theorem at a fixed error.
The stronger polynomial-error theorem is not required.
\end{abstract}

\section*{Proof status and contributions}
The complete assembly has separate line audits by the root, path-stability, and four-regular agents.
The random-host agent assembled it and independently reviewed the supporting proofs and the exposition refactor.
The new theorem has not been formalized in Lean or externally peer reviewed.
Publication priority has not been established.
The constants and rank thresholds are very large.

The proof builds on the existing research programme of the shared directory.
Fable 5.2 supplied the parallel-link lemma and the nested-chain and restoration machinery.
The present continuation adds the uniform affine-forest estimates, the iterative packing argument without an excess parameter, and the exact-half assembly.
The lower-bound construction is due to Horak and Kim~\cite{HK}.

\clearpage
\tableofcontents

\section{Conventions and three established inputs}
Rank means $\beta(G)=|E(G)|-|V(G)|+\comp(G)$.
The binary cycle space $\mathcal Z(G)$ consists of all even edge sets and has order $2^{\beta(G)}$.
Write $\xi_G$ for its uniform random element and $c(G)$ for the number of simple cycles.
Loops and parallel edges are allowed and remain distinct.
A loop is a cycle of length one, and a parallel pair is a cycle of length two.
Subdivision preserves rank and cycle count, so multigraphs and simple graphs have the same extremal value.

For a vertex set $P$, the notation $\xi_G-P$ means the restriction to the edges of $G-P$.
The notation $\xi_G\setminus E(W)$ instead removes only the specified edges.
These two operations are kept distinct throughout the proof.
A cyclic cut has a cycle on each shore.
A normalized $k$-pole is connected, has $k$ degree-two vertices called ports, and has degree three elsewhere.
All logarithms have base two, unless written as $\ln$.

The following three results were proved independently before the new assembly.
Their sources are specified so that the dependency chain can be checked.

\begin{theorem}[Established qualitative marked-cycle estimate]
\label{thm:sharp-qualitative-forest-input}
There is an absolute rank threshold $R_*$ such that the following holds.
Let $H$ be a connected simple graph of rank $\rho\ge R_*$.
If a marked cycle $W$ satisfies
\[
 |E(W)|\le \frac{\rho}{(\log\rho)^2},
\]
then
\[
 \Pr\bigl(\xi_H\setminus E(W)\text{ is a forest}\bigr)\le\frac{11}{20}.
\]
\end{theorem}

This is the one-cycle case of Theorem \texttt{thm:forest} and Remark \texttt{rem:range} in the original \nolinkurl{pancyclic-paper.tex}.
That theorem gives $1/2+20/\log^{(3)}\rho$, where $\log^{(3)}$ is the third iterated logarithm.
Choose a finite threshold with $\log^{(3)}\rho\ge400$ and $\log\rho\ge100$.
The remark permits the displayed marked-cycle size.
For the multigraph caps used below, subdividing only edges outside the marked cycle gives the same conclusion.
This preserves rank and the forest event, and leaves the marked cycle unchanged.

\begin{theorem}[Established sharp bound for small cyclic cuts]
\label{thm:sharp-order-four-cut-error}
There is an absolute constant $K_4$ such that every connected loopless cubic multigraph of rank $r$ with a cyclic cut of size at most four satisfies
\[
 c(G)\le2^{r-1}+K_4r2^{r/2}.
\]
\end{theorem}

The proof is the theorem ``Sharp-order four-cut error in \nolinkurl{small-cut-sharp-bound.pdf}.
It combines the four-pole profile classification with the linear two-cycle packing bound.
In the aligned saturated case, the largest cap product cancels in the exact gluing identity.
This input does not assume an unrestricted sharp-order estimate.

\begin{corollary}[Established deficit at bounded girth]
\label{cor:path-unbounded-girth-reduction}
For every fixed integer $g_0\ge2$, there is a rank threshold $R(g_0)$ such that the following holds.
A connected loopless cubic multigraph $G$ with rank at least $R(g_0)$, girth at most $g_0$, and no cyclic cut of size at most four satisfies
\[
 c(G)<2^{\beta(G)-1}.
\]
\end{corollary}

This is proved in \nolinkurl{path-fixed-girth.tex} and the separate bounded-girth and structural-reduction packets.
Its proof partitions terminal paths by their first departure from a shortest cycle.
After a prefix of $p$ vertices is removed, the target component has rank at most $r-1-p$.
The previously proved high-density path theorem and the asymptotic cycle bound yield, at each fixed girth $g\ge5$,
\[
 c(G)\le2^{r-1}-2^{r-g-1}+o_g(2^r).
\]
At sufficiently large rank, a cycle of length at most four itself gives a cyclic cut of size at most four.
Thus the displayed corollary covers every bounded girth.
Its proof does not use the new sharp-order conclusion.

The next two sections derive all new local inputs from the first established theorem.
The walk, link, chain, and assembly sections then prove the unrestricted result using the other two inputs.

\section{Qualitative affine-forest bounds for the sharp bootstrap}
\label{sec:path-qualitative-affine-bootstrap}

This section uses one qualitative marked-forest input.
It does not use a quantitative rate of convergence to one half.
All logarithms have base two.
A normalized $k$-pole is connected, has $k$ vertices of degree two, and has degree three at every other vertex.
Ranks are cycle-space dimensions, including for disconnected graphs.

\paragraph{Established forest input.}
Use Theorem~\ref{thm:sharp-qualitative-forest-input}, and increase its absolute rank threshold $R_*$ to at least $2^{100}$.
Its fixed-error conclusion is the following: if a connected graph has rank $\rho\ge R_*$ and a marked cycle $W$ with
\[
 |E(W)|\le\frac{\rho}{(\log\rho)^2},
\]
then
\begin{equation}\label{eq:path-qualitative-forest-input}
 \Pr(\xi\setminus E(W)\text{ is a forest})\le\frac{11}{20}.
\end{equation}
The input for simple graphs also applies to the localized multigraphs below.
Subdivide edges outside $W$ to remove loops and parallel edges; this preserves rank, the marked cycle, and the outside-forest probability.
Only this one-cycle input is used.

\begin{lemma}[Localization with a connected exterior]
\label{lem:path-qualitative-localization}
Let $G$ be connected, and let $R$ be a nonempty proper vertex set whose complement is connected.
For each edge of $\delta_G R$, add a distinct new vertex and an edge joining it to the original endpoint in $R$.
Join these new vertices in a cycle, padding to length three if necessary.
Call the new cycle $W_R$, its resulting graph $G_R$, and write $d=|\delta_G R|$.
Then $G_R$ is connected and
\[
 \beta(G_R)=\beta(G[R])+d-\operatorname{comp}(G[R])+1.
\]
Uniform even subgraphs have the same restriction law on $E(G[R])$ in $G$ and $G_R$.
If every vertex of $R$ is cubic in $G$, then
\[
 \beta(G_R)=\frac{|R|+d}{2}+1.
\]
\end{lemma}

\begin{proof}
Every component of $G[R]$ has a boundary edge, so the new cycle connects them all.
Before padding, the construction adds $d$ vertices and $2d$ edges.
Its rank is therefore $|E(G[R])|+d-|R|+1$, which is the first formula.
Padding adds equally many edges and vertices.
The cubic degree identity $3|R|=2|E(G[R])|+d$ gives the second formula.

Contract the connected complement of $R$ in $G$ and, separately, contract $W_R$ in $G_R$.
The resulting graphs agree on the original interior edges and the boundary edges.
Restriction under either contraction maps the original cycle space onto the quotient cycle space.
To see surjectivity, complete each prescribed even boundary by a join in a spanning tree of the contracted connected set.
A surjective linear map sends the uniform law to the uniform law on its image.
This proves the equality of restriction laws.
\end{proof}

\begin{lemma}[A fixed-error layer-width estimate]
\label{lem:path-qualitative-marker-width}
There is an absolute constant $B$ with the following property.
Let $G$ be connected and subcubic, and let $P$ be a nonempty connected marked set containing every vertex of degree below three.
Suppose that $P$ meets each distance layer from a fixed root in at most $k$ vertices, where $k\ge1$.
If
\[
 \Pr(\xi_G-P\text{ is a forest})\ge\frac35,
\]
then each full distance layer has at most $B(k+1)^2$ vertices.
\end{lemma}

\begin{proof}
Put $\Phi(m)=m/(\log m)^2$ and
\[
 \Theta=\max\{2R_*,\,2^{100}(k+1)^2\}.
\]
We first exclude every set $R\subseteq V(G)\setminus P$ with
\[
 m:=|R|\ge\Theta,
 \qquad |\delta_G R|\le\Phi(m)/4.
\]
Let $C_0$ be the component of $G-R$ containing $P$, and set
$R^+=V(G)\setminus V(C_0)$.
It contains $R$, avoids $P$, and has connected complement.
Distinct components of $G-R$ have no edge between them, so
$\delta_G R^+\subseteq\delta_G R$.
Every vertex of $R^+$ is cubic.

Localize at $R^+$ by Lemma~\ref{lem:path-qualitative-localization}.
Writing $d^+=|\delta_G R^+|$, the localized rank is
\[
 \rho=\frac{|R^+|+d^+}{2}+1\ge m/2\ge R_*.
\]
The function $\Phi$ is increasing on $[8,\infty)$.
If $d^+\ge3$, the marked cycle has order
\[
 d^+\le\Phi(m)/4\le\Phi(m/2)\le\Phi(\rho).
\]
If $d^+<3$, its padded order is three, which also satisfies this bound because $\rho\ge2^{100}$.
The event of a forest outside $P$ implies the forest event on $G[R^+]$.
In the localized graph, deleting the marked cycle edges leaves the new spokes incident with leaves.
Those spokes cannot create a cycle, so the outside-edge forest event is exactly the forest event on $G[R^+]$.
Localization and \eqref{eq:path-qualitative-forest-input} therefore bound its probability by $11/20$, contradicting the assumed lower bound $3/5$.
Thus every such $R$ is excluded.

Write $m_j$ for the number of unmarked vertices in distance layer $j$.
For the union $R(i,i')$ of unmarked layers $i$ through $i'$, its boundary satisfies
\[
 |\delta_G R(i,i')|
 \le3m_{i-1}+3m_{i'+1}+3k(i'-i+3).
\]
Indeed, every edge joins the same or consecutive layers; the last term counts incidences to marked vertices in layers $i-1$ through $i'+1$.
Use zero masses outside the layer range.

If some $m_j\ge\Theta$, begin with that layer.
Add a neighboring layer whenever its mass exceeds $m/(48(\log m)^2)$, where $m$ is the current mass.
The process stops in the finite graph.
Let $\mu$ be its final mass and $L=\log\mu\ge100$.
Each step multiplies the current mass by more than $1+1/(48L^2)$.
There are therefore at most
\[
 \frac{\ln\mu}{\ln(1+1/(48L^2))}
 \le96L^2\ln\mu<67L^3
\]
steps.
At termination the two neighboring masses are at most $\mu/(48L^2)$.
Hence
\[
 |\delta_G R(i,i')|
 \le\frac{\mu}{8L^2}+3k(67L^3+3)
 \le\frac{\mu}{4L^2}.
\]
For the last comparison, use
\[
 24kL^2(67L^3+3)\le2000kL^5\le2^L.
\]
The first inequality holds for $L\ge100$.
The threshold on $\mu$ gives $k\le2^{L/2-50}$, and
$2000L^5\le2^{L/2+50}$ for $L\ge100$.
The latter inequality holds at $100$, with an increasing right-to-left ratio thereafter.
The stopped window is thus an excluded set, a contradiction.

Every layer has fewer than $\Theta$ unmarked vertices and at most $k$ marked vertices.
For example,
\[
 B=2R_*+2^{100}+1
\]
bounds their total by $B(k+1)^2$.
\end{proof}

\begin{lemma}[Translation by a breadth-first tree marker]
\label{lem:path-qualitative-affine-width}
Let $N$ be a normalized $k$-pole with $k\ge2$.
If some even port subset $T$ has affine forest proportion at least $3/5$, then every distance layer from any chosen port has at most $B(k+1)^2$ vertices.
\end{lemma}

\begin{proof}
Choose a breadth-first tree rooted at that port and take the union of its root-to-port paths.
Its vertex set $P$ is connected, contains every degree-two vertex, and meets each distance layer in at most $k$ vertices.
The union of the tree edges contains a $T$-join $F_0$.
Translation by $F_0$ is a bijection from the even space to the affine $T$-join space.
The translated and original states agree outside $P$.
Thus an affine forest gives a state that is a forest outside $P$, and the preceding lemma applies.
\end{proof}

\begin{lemma}[Weighted Moore bound after suppressing the ports]
\label{lem:path-qualitative-weighted-moore}
Let $N$ be a normalized $k$-pole of rank $h\ge2$ and girth at least $D$.
With $t=\lceil\log(8h)\rceil$, one has
\[
 D\le4t\left(1+\frac{k}{3(h-1)}\right).
\]
\end{lemma}

\begin{proof}
Suppress all degree-two vertices and give each resulting edge the length of the path that it replaces.
The kernel is a connected cubic multigraph with
$n_0=2h-2$ vertices and $m_0=3h-3$ edges.
Its positive edge lengths sum to $m_0+k$.
Loops and parallel edges are permitted, and weighted cycles lift to cycles of the same length in $N$.

The kernel has $2m_0\,2^{t-1}$ directed non-backtracking walks of $t$ edges.
Under the uniform law on these walks, each directed edge is uniform at every step.
The expected total weight is $t(1+k/m_0)$.
At least half the walks have weight at most twice that expectation, and their number is at least
\[
 m_0\,2^{t-1}\ge12h(h-1)>n_0^2.
\]
If a short walk repeats a vertex, its support contains a cycle of no greater weight.
Otherwise all these short walks are simple paths, and two have the same ordered endpoints.
Their union contains a cycle of weight at most the sum of their weights.
In either case, a cycle has weight at most $4t(1+k/m_0)$.
This proves the bound.
\end{proof}

\begin{lemma}[The fixed-ratio affine input]
\label{lem:path-qualitative-fixed-ratio-core}
There is an absolute integer $D_*\ge3$ such that every normalized $k$-pole with
\[
 h\ge256,\qquad k\le4D,
 \qquad\operatorname{girth}(N)\ge D\ge D_*
\]
has affine forest proportion at most $3/5$ for every even port subset.
The assertion includes $k=0,1$.
\end{lemma}

\begin{proof}
With at most one port, the only even port subset is empty, and its forest proportion is $2^{-h}$.
Assume $k\ge2$ and put $t=\lceil\log(8h)\rceil$.
For $h\ge256$,
\[
 h-1>\frac{32}{3}\,t.
\]
Indeed, $t\le\log h+4$, the ratio $(h-1)/(\log h+4)$ increases for $h\ge2$, and at $h=256$ it exceeds $32/3$.
The weighted Moore bound therefore gives
\[
 D\le4t+\frac{16tD}{3(h-1)}<4t+\frac D2.
\]
Thus
\[
 h\ge2^{D/8-4},\qquad n=|V(N)|\ge h.
\]

Put $s=\lfloor D/16\rfloor$.
A graph of maximum degree three has at most $3\cdot2^s-2$ vertices in each radius-$s$ ball.
The union of these balls around the ports has at most
\[
 3k2^s\le12D2^{D/16}<2^{D/8-4}\le n
\]
vertices for sufficiently large $D$.
Choose a vertex outside that union.
Its radius-$s$ ball contains no port.
Since $2s<D$, its cubic breadth-first expansion has at least
$2^s\ge2^{D/16-1}$ vertices.
This ball meets at most $2s+1\le D$ distance layers from any fixed root.
Some layer therefore has at least $2^{D/16-1}/D$ vertices.

If any affine forest proportion were greater than $3/5$, Lemma~\ref{lem:path-qualitative-affine-width} would bound that layer by
$B(k+1)^2\le B(4D+1)^2$.
Choose $D_*$ large enough that the preceding ball comparisons hold and
\[
 \frac{2^{D/16-1}}D>B(4D+1)^2
 \qquad(D\ge D_*).
\]
The exponential term makes such a fixed threshold possible.
This contradiction proves the lemma.
\end{proof}

\begin{lemma}[An affine bound at the boundary scale]
\label{lem:path-clean-core-affine-forest}
There is an absolute constant $C_0$ such that every normalized $k$-pole with
\[
 h\ge L\ge4,
 \qquad k\le2L,
 \qquad\operatorname{girth}(N)\ge D\ge C_0\log(L+2)
\]
has affine forest proportion at most $3/5$ for every even port subset.
The assertion includes $k=0,1$.
\end{lemma}

\begin{proof}
The cases $k\le1$ again have forest proportion $2^{-h}$.
For $k\ge2$, one has
\[
 1+\frac{k}{3(h-1)}\le1+\frac{2h}{3(h-1)}<2.
\]
The weighted Moore bound gives
\[
 D<8\lceil\log(8h)\rceil\le8\log h+32,
 \qquad h\ge2^{D/8-4}.
\]
Choose $s=\lfloor D/16\rfloor$.
The radius-$s$ port balls contain at most $6L2^{D/16}$ vertices.
For a sufficiently large absolute $C_0$, this is less than
$n\ge h\ge2^{D/8-4}$ whenever $D\ge C_0\log(L+2)$.
The resulting port-free cubic ball gives a layer of size at least
$2^{D/16-1}/D$.

An affine forest proportion greater than $3/5$ would instead bound the layer by
$B(k+1)^2\le B(2L+1)^2$.
Increase $C_0$ so that the lower bound exceeds this quantity uniformly for $L\ge4$.
For explicit uniformity, the function $2^{D/16}/D$ increases for $D\ge32$.
At $D=C_0\log(L+2)$, its ratio to $(L+2)^2$ is
\[
 \frac{(L+2)^{C_0/16-2}}{C_0\log(L+2)}.
\]
This tends uniformly to infinity as $C_0\to\infty$, for $L\ge4$.
The same choice ensures the port-ball comparison.
The resulting contradiction proves the lemma.
\end{proof}

\begin{lemma}[Extracting a core from a large block]
\label{lem:path-clean-block-core}
Let $G$ be connected and cubic, and let $R$ be a nonempty proper vertex set with
\[
 |R|\ge8L^2,\qquad |\delta_G R|\le2L,\qquad L\ge4.
\]
Some connected component of $G[R]$ has a connected two-core of rank at least $L$, with at most $2L$ degree-two vertices and degree three elsewhere.
\end{lemma}

\begin{proof}
Write $n=|R|$, $d=|\delta_G R|$, and $c=\operatorname{comp}(G[R])$.
Each component has a boundary edge, so $1\le c\le d\le2L$.
The cubic degree identity gives
\[
 \beta(G[R])=\frac{n-d}{2}+c.
\]
Thus the average component rank is at least
\[
 \frac{n-d}{2d}+1
 \ge\frac{8L^2-2L}{4L}+1>L.
\]
Take a component of rank at least $L$ and prune its degree-one vertices.
The nonempty two-core remains connected and keeps the same rank.
At each step, the removed leaf has two incidences outside the current set and one inside it, so the boundary decreases by one.
The core boundary is therefore at most $2L$.
Its vertices have internal degree two or three, and each degree-two vertex has one boundary incidence.
This proves the port bound.
\end{proof}

\begin{proposition}[Two clean blocks]
\label{prop:path-two-clean-blocks}
Let $G$ be connected and cubic, and let $P$ be any marked set.
Suppose two vertex-disjoint sets $R_1,R_2\subseteq V(G)\setminus P$ satisfy
\[
 |R_i|\ge8L^2,\qquad |\delta_G R_i|\le2L,
 \qquad\operatorname{girth}(G[R_i])\ge D,
\]
where $L\ge4$ and $D\ge C_0\log(L+2)$.
Then
\[
 \Pr(\xi_G-P\text{ is a forest})\le\frac9{25}<\frac12.
\]
\end{proposition}

\begin{proof}
Extract cores $K_1,K_2$ by Lemma~\ref{lem:path-clean-block-core}.
Each remains disjoint from $P$, has rank at least $L$, at most $2L$ ports, and girth at least $D$.
Condition on every edge outside $E(K_1)\cup E(K_2)$.
For each feasible conditioning, the remaining states are independent and uniform in two affine parity classes.
The fiber is the product of their internal cycle spaces, translated by any compatible pair of states.
Only degree-two vertices can have odd prescribed parity, and each prescribed odd set has even size.

Lemma~\ref{lem:path-clean-core-affine-forest} bounds each conditional forest probability by $3/5$.
The forest event outside $P$ requires both core states to be forests.
Its conditional probability is therefore at most $9/25$.
Average over the conditioned edges.
Edges joining the two cores are already conditioned, so nonadjacency is not needed.
\end{proof}

\section{Cycle packing with no positive-excess parameter}
\label{sec:four-unrestricted-iterative-packing}

Throughout this section, $\xi_G$ is a uniform member of the binary cycle space of $G$.
The rank of a graph is its cycle-space dimension.
All logarithms have base two.
We use only the following qualitative input from
Lemma~\ref{lem:path-qualitative-fixed-ratio-core}:
there is an absolute $D_*\ge3$ such that, with $H=256$, every connected normalized pole of rank at least $H$, girth at least $D\ge D_*$, and at most $4D$ degree-two ports has forest probability at most $3/5$ in every feasible affine parity class.
All its other vertices have degree three.

\begin{lemma}[Terminal bipartitions along a forest]
\label{lem:four-terminal-bipartition}
Let $Q$ be a connected loopless multigraph and let $A$ be a nonempty proper vertex set.
For each vertex $v$, let $L_v$ be the span of its incident coordinate forms on $\mathcal Z(Q)$.
Put $U=\sum_{a\in A}L_a$ and $B=|\delta_QA|$.
Among vertices $v\notin A$ with $L_v\cap U\ne0$, at most $2B$ are exceptional.
The others form at most $2B$ ordered classes.
Within each class, the vertices are articulation vertices of one component of $Q-A$ and give the same bipartition of its boundary terminals.
Each edge to $A$ supplies a distinct terminal.
\end{lemma}

\begin{proof}
Exclude the at most $B$ vertices adjacent to $A$.
In each component of $Q-A$, choose a spanning tree, attach a leaf for each boundary terminal, and retain the union of paths between those leaves.
This forest has at most $B$ leaves and at most $B$ branching vertices.
Exclude its branching vertices and suppress its degree-two vertices.
The resulting forest has at most $2B$ edges.

We verify that those edges contain all other relevant vertices.
Fix such a vertex $v$ and choose $0\ne a\in L_v\cap U$.
Represent $a$ both by a subset of the star at $v$ and by edges incident with $A$.
Their symmetric difference annihilates the cycle space, so it is a cut.
This elementary cut-space fact follows, for example, by assigning each vertex the parity of a root-to-vertex tree path; orthogonality to every fundamental cycle gives the same endpoint rule on each nontree edge.

Let $K$ be the component of $Q-A$ containing $v$.
Every component of $K-v$ lies entirely on one side of that cut.
Thus the star subset joins $v$ to a union of components of $K-v$.
As a functional, it equals the sum of the terminal coordinates in that union.
The sum of all terminal coordinates of $K$ is zero.
Since $a\ne0$, at least two components of $K-v$ contain terminals.
Every spanning tree must connect those terminal groups through $v$.
After removing the branching exceptions, exactly two terminal-bearing components remain.
Their terminal sets are the two sides of one edge of the suppressed forest.
Vertices on the same edge therefore give the same actual terminal bipartition.
Ordering them along that edge proves the assertion.
\end{proof}

\begin{lemma}[Two clean regions]
\label{lem:four-two-clean-regions}
Let $G$ be connected and cubic, let $P$ be connected or empty, and let $S$ be disjoint from $P$.
Suppose $G-(P\cup S)$ has girth at least $D\ge D_*$.
If $R_1,R_2$ are disjoint connected vertex sets avoiding $P\cup S$, each of rank at least $H$ and boundary at most $4D$, then
\[
 \Pr(\xi_G-(P\cup S)\text{ is a forest})\le\frac9{25}.
\]
\end{lemma}

\begin{proof}
Prune leaves from each induced graph $G[R_i]$.
Its connected two-core has the same rank and no larger boundary.
Indeed, removing a leaf in a cubic ambient graph replaces two external incidences by one, decreasing the boundary by one.
Each core has internal degrees two and three, with one boundary edge at each degree-two vertex.
It retains girth at least $D$.

Condition on every edge outside the two core interiors, including any edges between the cores.
For every conditioning of positive probability, the internal fiber is the Cartesian product of two affine translates of the internal cycle spaces.
Thus the two internal states are independent and uniform in their parity classes.
Odd prescribed parity can occur only at degree-two ports, and its port set has even size.
Each conditional forest probability is at most $3/5$ by the qualitative core lemma.
If a core has zero or one port, its only feasible parity class is the even space, whose only forest is empty and has probability $2^{-h}\le3/5$.
The global forest event requires both restrictions to be forests, so its probability is at most $(3/5)^2$.
\end{proof}

\begin{proposition}[Packing after a previous deletion]
\label{prop:four-packing-after-deletion}
Under the preceding notation, suppose
\[
 \Pr(\xi_G-(P\cup S)\text{ is a forest})>\frac12,
 \qquad s=|S|.
\]
Let $2\le d\le2D$.
Every family of $q$ disjoint connected cyclic regions in $G-(P\cup S)$, each with boundary in $G$ at most $d$, satisfies
\[
 q<(6H+15)d(s+1)2^d.
\]
\end{proposition}

\begin{proof}
If $q=0$ there is nothing to prove.
Contract each region to a marked vertex, and contract each connected component outside all regions to one unmarked vertex.
Discard loops, and call the resulting connected quotient $Q$.
Its marked vertices have degree at most $d$.

Restriction to the quotient edges maps $\mathcal Z(G)$ onto $\mathcal Z(Q)$.
To see surjectivity, extend the even quotient boundary pattern inside each connected contracted part by any parity join.
The fiber is a product of affine internal cycle spaces, so the quotient state is uniform and its internal choices are conditionally independent.
At a marked vertex $v$, let $I_v$ indicate that every incident quotient edge is absent.
In that event the internal state is uniform even; since the region is cyclic, being a forest requires the empty state and has probability at most $1/2$.
When $I_v=0$, use the upper bound one.
It follows that
\[
 \Pr(\xi_G-(P\cup S)\text{ is a forest})
 \le\mathbb E\prod_v(1-I_v/2).
\]

Write $L_v$ for the incident coordinate span on $\mathcal Z(Q)$.
It has dimension at most $d-1$, because even star parity gives one relation.
No $L_v$ is zero: a zero space would make $I_v=1$ identically and the displayed bound at most one half.
Take a maximal family of these spaces forming a direct sum, with vertex set $A$ and size $t$.
For this family the zero-star events are independent, each with probability
$2^{-\dim L_v}\ge2^{1-d}$.
Discarding the other factors gives
\[
 \Pr(\xi_G-(P\cup S)\text{ is a forest})
 \le(1-2^{-d})^t.
\]
If $t\ge2^d$, this is at most $e^{-1}<1/2$.
Consequently $1\le t<2^d$.
If all marked vertices are selected, the required conclusion follows already.

Otherwise put $U=\sum_{a\in A}L_a$ and $B=|\delta_QA|\le dt$.
Maximality gives $L_v\cap U\ne0$ at every remaining marked vertex.
Lemma~\ref{lem:four-terminal-bipartition} places all but at most $2B$ of them in at most $2B$ ordered classes.

Consider one class $v_1,\ldots,v_m$.
Every consecutive block of $H+2$ class vertices gives a middle component of $Q-A$ after deleting its first and last vertices.
It contains the $H$ interior marked vertices.
It has no neighbor in $A$: any attached terminal in that middle component would change the actual terminal bipartition between the first and last vertices.
Every edge leaving it therefore ends at those two vertices, giving boundary at most $2d$.
The articulation separations make the middle components from disjoint consecutive blocks pairwise disjoint.
Their lifts to $G$ are connected and contain $H$ disjoint cyclic regions, hence have rank at least $H$.

If $P$ is nonempty, its connectedness and avoidance of all marked regions put it in one contracted outside component.
At most one middle lift can meet it.
At most $s$ further lifts meet $S$.
Thus $m\ge(H+2)(s+3)$ would give at least two clean lifts, each with boundary at most $2d\le4D$.
Lemma~\ref{lem:four-two-clean-regions} would make the forest probability at most $9/25$, a contradiction.
Every class therefore has fewer than $(H+2)(s+3)$ members.

Counting selected vertices, exceptions, and classes now gives
\begin{align*}
 q&\le t+2B+2B(H+2)(s+3)\\
  &\le t+(6H+14)(s+1)B\\
  &\le(6H+15)d(s+1)t
   <(6H+15)d(s+1)2^d.
\end{align*}
\end{proof}

\begin{theorem}[An excess-free girth-raising bound]
\label{thm:four-excess-free-girth-raising}
There is an absolute constant $K$ with the following property.
Let $G$ be connected and cubic and let $P$ be connected or empty.
Assume
\[
 \Pr(\xi_G-P\text{ is a forest})>\frac12,
 \qquad \operatorname{girth}(G-P)\ge D_*.
\]
For each integer $D\ge D_*$, there is $S\subseteq V(G)\setminus P$ such that
\[
 \operatorname{girth}(G-(P\cup S))>D,
 \qquad |S|+1\le2^{\,2D+K\log^2(D+2)}.
\]
The set $S$ is the vertex union of a disjoint cycle family, maximal among cycles of length at most $D$ outside $P$.
On the common forward dyadic scales $D_j=2^jD_*$, the sets $S_j$ can be chosen nested with this bound simultaneously at every scale.
\end{theorem}

\begin{proof}
Deletion can only increase the probability of the forest event.
To raise an available girth scale $a\ge D_*$ to $b\le2a$, take a maximal disjoint family of cycles of length at most $b$ outside $P\cup S$ and add its vertices to $S$.
Every cycle has boundary at most its length in the cubic ambient graph.
The packing proposition, applied with $d=b$, gives
\[
 s_{\mathrm{new}}+1
 \le(6H+16)b^2 2^b(s+1).
\]
This also holds for an empty family.
The remaining graph has girth greater than $b$.

For an arbitrary target $D$, plan the scales backwards as $\lceil D/2^j\rceil$, stopping at the first value between $D_*$ and $2D_*$.
Read them forwards, starting with $S=\varnothing$ and available girth $D_*$.
If $D=D_*$, perform one deletion at that same scale.
Every step at most doubles the available scale; the sum of the new scales is at most $2D+O(\log(D+2))$.
There are $O(\log(D+2))$ steps.
Taking logarithms of the recurrence gives the asserted bound after increasing the absolute constant $K$.

Alternatively, use the common forward sequence $D_*,2D_*,4D_*,\ldots$.
The scales through $D_j$ sum to $2D_j-D_*$, so the same estimate holds simultaneously, and deletions are nested.
At every stage new cycles avoid all earlier deleted vertices.
Thus the combined family is disjoint, and the final girth assertion is exactly its maximality at the target length.
\end{proof}

\paragraph{Verification.}
This fragment restates the independently audited packing argument without its earlier forbidden-cut formulation.
An exact checker verifies the terminal partitions, middle intervals, core pruning, and complete conditional affine products:
\begin{center}
\small\nolinkurl{four-high-connectivity-iterative-packing-check.py}.
\end{center}
It does not numerically test the absolute high-girth threshold.

% Extracted from the audited polynomial foundations; original source unchanged.
\section{Cycle-space, walk, and spectral estimates}
\label{sec:polynomial-foundations}

A region is a nonempty connected vertex set.
For a region $U$, let $\delta U$ be its edge boundary.
Write $\mathcal O_U$ for the event that a uniform even subgraph avoids $\delta U$.
For disjoint regions $U,V$, their links are the edges between them and the components of $G-U-V$ adjacent to both.
The number of links is $\ell(U,V)$.
The following proofs are adapted from Section6 of the existing manuscript.
They are included to expose every dependency of the sharp-order argument.

\begin{lemma}[Link formula]\label{lem:link}
Let $A$ and $B$ be disjoint regions of a connected graph $G$. Then
$\PP(\cO_A)=2^{\comp(G-A)-|\delta A|}\ge2^{-|\delta A|}$ and
\[
  \PP(\cO_A\cap\cO_B)=2^{\ell(A,B)-1}\,\PP(\cO_A)\,\PP(\cO_B).
\]
Conditionally on $\cO_A$, the set $\xi\cap E(G[A])$ is uniform on $\Z(G[A])$; conditionally on
$\cO_A\cap\cO_B$, the sets $\xi\cap E(G[A])$ and $\xi\cap E(G[B])$ are independent and uniform on
$\Z(G[A])$ and $\Z(G[B])$.
\end{lemma}

\begin{proof}
The even subgraphs that avoid $\delta A$ are the sums of an even subgraph of $G[A]$ and an even
subgraph of $G-A$. So $\PP(\cO_A)=2^{\beta(G[A])+\beta(G-A)-\beta(G)}$, and counting edges and
vertices shows that the exponent is $\comp(G-A)-|\delta A|$. In the same way
\[
  \PP(\cO_A\cap\cO_B)=2^{\beta(G[A])+\beta(G[B])+\beta(G-A-B)-\beta(G)}
     =2^{\comp(G-A-B)+1-|\delta A|-|\delta B|+e(A,B)},
\]
where $e(A,B)$ is the number of edges between $A$ and $B$. Let $a$, $b$ and $k$ be the numbers of
components of $G-A-B$ that are joined to $A$ only, to $B$ only, and to both. Then
$\comp(G-A-B)=a+b+k$, while $\comp(G-A)=a+1$ and $\comp(G-B)=b+1$, and the formula follows since
$\ell(A,B)=k+e(A,B)$. The conditional laws are read off from the two direct sum decompositions.
\end{proof}

So two cut events are independent when there is one link. They are correlated by a factor of
exactly two when there are two links. The correlation is stronger only for pairs with at least
three links, and the next lemma says that such pairs are rare.

\begin{lemma}[Exceptional partners]\label{lem:exc}
Let $A$ be a region of a connected graph $G$, and let $B_1,\dots,B_k$ be pairwise disjoint regions,
disjoint from $A$, with $\ell(A,B_i)\ge3$ for all $i$. Then $k\le|\delta A|$.
\end{lemma}

\begin{proof}
Contract $A$ to a vertex $a$ and each $B_i$ to a vertex $b_i$, and delete the loops that arise.
This changes neither $d=|\delta A|$ nor the numbers $\ell(A,B_i)$, because the sets $B_j$ with
$j\ne i$ are connected and are contracted inside components of $G-(A\cup B_i)$. Let $F$ be a spanning forest of
$G-a$, and for every edge $e=au$ at $a$ attach a new leaf $\tau_e$ to $u$. Call the resulting forest
$F^{+}$ and its $d$ new leaves \emph{terminals}.

Fix $i$. Every link between $a$ and $b_i$ yields a path of $F^{+}$ from $b_i$ to a terminal. For an
edge $e=ab_i$ it is the edge $b_i\tau_e$. For a component $Y$ of $G-\{a,b_i\}$ joined to both, choose
an edge $e=au$ with $u\in Y$ and take the path of $F^{+}$ from $b_i$ to $\tau_e$; after its first
vertex it stays in $Y\cup\{\tau_e\}$. Paths that come from different links therefore start with
different edges at $b_i$, and the union of any two of them is a path between two terminals. Let
$F^{\circ}\subseteq F^{+}$ be the union of all paths between terminals. As $\ell(a,b_i)\ge3$, the
vertex $b_i$ has degree at least $3$ in $F^{\circ}$. All leaves of $F^{\circ}$ are terminals, and a
forest with at most $d$ leaves has at most $d-2$ vertices of degree at least $3$. Hence $k\le d-2$ if $k\ge1$.
\end{proof}


\subsection{Walks and normalized traces}
Let $J$ be a simple graph of maximum degree three and girth greater than $D\ge3$.
Put $s=\lfloor(D-1)/2\rfloor$ and let $L>D$.
Let $N_\ell(u,w)$ count non-backtracking walks of length $\ell$ from $u$ to $w$.
Write $v=|V(J)|$ and $d_J=\sum_u(3-\deg_Ju)$.
For the spectral supply estimate below, assume $d_J/v<1$.

\begin{lemma}\label{lem:walks}
If $\ell>s$, then $N_\ell(u,w)\le3\cdot2^{\ell-s-1}$ for all $u,w$. Consequently, for every vertex
$u$ of $J$, the cycles $C$ of $J$ through $u$ with $|C|\le L$ satisfy $\sum_C2^{-|C|}\le L\,2^{-s}$.
\end{lemma}

\begin{proof}
There are at most $3\cdot2^{\ell-s-1}$ ways to choose the first $\ell-s$ steps. The last $s$ steps
form a non-backtracking walk of length $s$ between two given vertices. Such a walk is a path,
because $s$ is less than the girth, and there is at most one such path, because two of them would
give a cycle of length at most $2s<D$. For the second statement, a cycle of length $\ell$ through $u$
can be traversed from $u$ in two directions, so there are at most $\half N_\ell(u,u)$ of them, and
$\ell>D>s$.
\end{proof}

Let $B$ be the non-backtracking matrix of $J$. Its rows and columns are indexed by the directed
edges $(i\to j)$ of $J$, two for each edge. The entry in row $(i\to j)$ and column $(j'\to k)$ is
$1$ if $j'=j$ and $k\ne i$, and it is $0$ otherwise.

\begin{lemma}\label{lem:spectrum}
For every even $\ell\ge2$,
\[
  \tr B^{\ell}\;\ge\;\bigl(2-d_J/v\bigr)^{\ell}-3v\,2^{\ell/2}.
\]
\end{lemma}

\begin{proof}
Let $A$ be the adjacency matrix of $J$ and $\Delta$ the diagonal matrix of its degrees, and put
$M(\lambda)=\lambda^{2}I-\lambda A+\Delta-I$.

(1) \emph{If $M(\lambda)g=0$ for a real $\lambda>1$ and a vector $g\ne0$, then $\lambda$ is an
eigenvalue of $B$.} A direct computation shows that $f(i\to j)=\lambda g(j)-g(i)$ satisfies
$Bf=\lambda f$; and $f=0$ would give $g(i)=\lambda g(j)=\lambda^{2}g(i)$ along every edge, hence
$g=0$ (at an isolated vertex, $M(\lambda)g=0$ already says $(\lambda^{2}-1)g=0$).

(2) \emph{If $Bf=\lambda f$ with $f\ne0$ and $\lambda^{2}\ne1$, then $M(\lambda)g=0$ for
$g(j)=\sum_{k\sim j}f(j\to k)$, and $g\ne0$.} The eigenvalue equation at $(i\to j)$ and at
$(j\to i)$ reads $\lambda f(i\to j)=g(j)-f(j\to i)$ and $\lambda f(j\to i)=g(i)-f(i\to j)$.
Eliminating $f(j\to i)$ gives $(\lambda^{2}-1)f(i\to j)=\lambda g(j)-g(i)$. This shows that
$g\ne0$. Summing it over the neighbours $j$ of $i$ gives $M(\lambda)g=0$.

(3) \emph{$B$ has a real eigenvalue $\lambda_1\ge2-d_J/v$.} The matrix
$M(2)=2(\Delta-A)+(3I-\Delta)$ is positive semidefinite, while
$\mathbf{1}^{\!\top}M(\lambda)\mathbf{1}=(\lambda-1)(v\lambda-2v+d_J)$ vanishes at
$\lambda^{*}=2-d_J/v>1$. The smallest eigenvalue of the symmetric matrix $M(\lambda)$ is continuous
in $\lambda$, at most $0$ at $\lambda^{*}$ and at least $0$ at $2$. So $M(\lambda_1)$ is singular for
some $\lambda_1\in[\lambda^{*},2]$, and (1) applies.

(4) \emph{Every non-real eigenvalue $\lambda$ of $B$ satisfies $|\lambda|\le\sqrt2$.} With $g$ as
in~(2), $g^{*}M(\lambda)g=0$ is a quadratic equation $a\lambda^{2}-b\lambda+c'=0$ with real
coefficients $a=g^{*}g>0$, $b=g^{*}Ag$ and $c'=g^{*}(\Delta-I)g\le2a$. Its roots are $\lambda$ and
$\bar\lambda$, so $|\lambda|^{2}=c'/a\le2$.

Now $\tr B^{\ell}$ is the sum of the $\ell$th powers of the $2|E(J)|\le3v$ eigenvalues of $B$. For
even $\ell$ the real eigenvalues contribute nonnegative terms, one of which is at least
$(2-d_J/v)^{\ell}$, and the non-real ones contribute at least $-3v\,2^{\ell/2}$ in total.
\end{proof}

The argument uses the same matrix-pencil proof as the existing manuscript.
It needs no expansion assumption.


\begin{lemma}[Repeated closed walks]\label{lem:polynomial-repeated-walks}
Let $\mathcal C$ contain the cycles of $J$ with length at most $L$ and put $\mu=\sum_{C\in\mathcal C}2^{-|C|}$.
Then
\[
 \mu\le T:=\sum_{\ell\le L}\frac{\operatorname{tr}B^\ell}{2\ell2^\ell}
 \le\left(1+\tfrac32L^2 2^{-s}\right)\mu.
\]
\end{lemma}
\begin{proof}
$\tr B^{\ell}$ is the number of closed walks $u_0u_1\cdots u_\ell=u_0$ in $J$ with
$u_{i+1}\ne u_{i-1}$ for all $i$ modulo $\ell$. A cycle of length $\ell$ gives $2\ell$ of them; let
$n_\ell$ be the number of the others. Then
\[
  T:=\sum_{\ell\le L}\frac{\tr B^{\ell}}{2\ell\,2^{\ell}}=\mu+\sum_{\ell\le L}\frac{n_\ell}{2\ell\,2^{\ell}} .
\]
Consider a walk counted by $n_\ell$, and let $j$ be the least index with
$u_j\in\{u_0,\dots,u_{j-1}\}$, say $u_j=u_i$. Then $j<\ell$, the walk $C=u_iu_{i+1}\cdots u_j$ is a
cycle, of length $q$ say, and $u_ju_{j+1}\cdots u_\ell u_1\cdots u_i$ is a non-backtracking closed
walk of length $\ell-q>0$ from $u_i$ to itself. It contains a cycle, so $\ell-q>D>s$. The original
walk is determined by $C$, by the vertex $u_i$ and a direction on $C$ ($2q$ choices), by the
complementary walk (at most $3\cdot2^{\ell-q-s-1}$ choices, by Lemma~\ref{lem:walks}), and by the
position of $u_0$ on the latter (at most $\ell$ choices). Hence
\[
  \sum_{\ell\le L}\frac{n_\ell}{2\ell\,2^{\ell}}
  \;\le\;\sum_{\ell\le L}\ \sum_{C\in\mathcal{C}}\tfrac32\,|C|\,2^{-|C|-s}
  \;\le\;\tfrac32\,L^{2}\,2^{-s}\,\mu,
\]
as required.


\end{proof}

In particular, suppose $v=O(r)$, $x=\log_2r$, $L=\lfloor Ax\rfloor$, and $A>6$ is fixed.
If $Ld_J=o(v)$ and $L^2 2^{-s}=o(1)$, then
\[
 \mu\ge(1/4-o(1))\ln(A/6).
\]
Indeed, sum the even trace bound over $6x\le\ell\le L$.
Its normalized leading term is $1-o(1)$, and its nonreal error is $O(r^{-2})$.
The even harmonic sum is $(1/4-o(1))\ln(A/6)$.
The repeated-walk error is $o(\mu)$.

\subsection{Component variables and exceptional pairs}
Let $\Gamma$ be the ambient connected cubic graph, and let $J$ be its induced unmarked subgraph with the preceding girth and degree bounds.
The uniform even state below is sampled from $\mathcal Z(\Gamma)$.
For a cycle $C\subseteq J$, put $a_C=2^{-|C|}$ and let $I_C$ mean that $C$ is a component of the uniform even subgraph $\xi$ of $\Gamma$.
Set $p_C=\Pr(I_C)$ and $X_C=(a_C/p_C)\mathbf1_{I_C}$.
The condition $E(C)\subseteq\xi$ is feasible and imposes at most $|C|$ binary linear equations.
Hence $p_C\ge a_C$, so $\mathbb EX_C=a_C$ and $\mathbb EX_C^2\le a_C$.
The link formula gives
\[
 \mathbb E(X_CX_{C'})=a_Ca_{C'}2^{\ell(C,C')-1}
\]
for disjoint cycles.
Distinct intersecting cycles cannot both be components.

Fix $C$ and choose a maximal disjoint family of partners with at least three links to $C$.
The exceptional-partner lemma bounds its size by $L$.
Its vertex union has size at most $L^2$ and meets every exceptional partner.
The local walk bound therefore bounds their total weight by $L^3 2^{-s}$.
If all link counts are at most $b$, their total normalized joint contribution is at most
\[
 a_C2^bL^3 2^{-s}.
\]
Only the larger girth scale will need this exceptional-pair estimate.

\section{Parallel links, the exact-half moment, and a nested chain}
\label{sec:sharp-parallel-links}

The parallel-link argument is due to Fable 5.2.
Its bound depends only on the strict inequality above one half.
No positive lower bound on the excess is needed.

\begin{lemma}[Fable's parallel-link estimate]
\label{lem:endpoint-parallel-links}
Let $A,B$ be disjoint connected vertex sets in a connected multigraph $\Gamma$.
Let $P$ be a vertex set, and let $\xi$ be a uniform even subgraph.
Suppose its restriction outside $P$ is a forest with probability greater than one half.
For every integer $t\ge0$, at most $2^{t+1}$ components of $\Gamma-(A\cup B)$ satisfy all four conditions below:
\begin{enumerate}
\item They have an edge to each of $A,B$.
\item They avoid $P$.
\item They contain a cycle.
\item Their boundary has at most $t$ edges.
\end{enumerate}
\end{lemma}

\begin{proof}
If there is no component joining $A$ to $B$, the assertion is immediate.
Otherwise retain one such component $Q_0$.
The set $A\cup B\cup Q_0$ is connected.
For any two other complementary components $Q_i,Q_j$, the graph $\Gamma-Q_i-Q_j$ is connected.
It contains the retained connected set, and every other complementary component has an edge to that set.
No edge joins $Q_i$ and $Q_j$.

Thus the two regions have exactly one link.
The cycle-space pair law makes their zero-boundary events independent.
Conditional on both zero-boundary events, their internal even edge sets are independent and uniform.
This also follows from the direct product of the cycle spaces after deleting the two boundaries.

For each eligible component other than $Q_0$, let $E_i$ mean that its boundary is zero and its internal even edge set is nonzero.
The events $E_i$ are pairwise independent.
The zero-boundary condition imposes at most $t$ linear equations, so its probability is at least $2^{-t}$.
The cyclic internal graph has a nonzero even edge set with conditional probability at least one half.
Consequently $p_i:=\Pr(E_i)\ge2^{-t-1}$.

If there are $m$ such components, put $N=\sum_i\mathbf1_{E_i}$ and $\nu=\mathbb EN$.
Pairwise independence gives
\[
 \mathbb EN^2\le\nu^2+\nu,\qquad
 \Pr(N=0)\le\frac1{1+\nu},\qquad
 \nu\ge m2^{-t-1}.
\]
The forest event implies $N=0$.
Its probability exceeds one half, so $m<2^{t+1}$.
Including $Q_0$, if it was eligible, gives at most $2^{t+1}$ components.
\end{proof}


\begin{lemma}[A uniform link bound]
\label{lem:sharp-uniform-links}
Let $G$ be connected and cubic, and let $P$ be a nonempty connected marked set.
Assume that a uniform even subgraph is a forest outside $P$ with probability greater than one half.
Let $S$ be the union of a maximal vertex-disjoint family of cycles of length at most $D$ outside $P$.
Put $J=G-(P\cup S)$.
If $8\le D\le L$, $t\ge1$, and $2(t+2\lfloor D/8\rfloor)\le D$, then disjoint cycles $C,C'$ of $J$, each of length at most $L$, satisfy
\[
 \ell(C,C')\le1+\frac{2L}{t+1}+2^{t+1}+578(L/D)^2.
\]
\end{lemma}

\begin{proof}
Put $U=V(C)\cup V(C')$.
At most one component of $G-U$ meets $P$, because $P$ is connected.
At most $2L/(t+1)$ components have more than $t$ boundary edges: the total boundary of $U$ is at most $2L$.
Among the other components joining the two cycles, at most $2^{t+1}$ are cyclic by Lemma~\ref{lem:endpoint-parallel-links}.

Every remaining component $Y$ that joins both cycles is a tree, avoids $P$, and has boundary at most $t$.
The cubic degree identity gives $|Y|=|\delta Y|-2\le t-2$.
It contains no vertex of $S$.
Indeed, each chosen cycle avoids $U$, so any chosen cycle meeting $Y$ would lie entirely in that tree.
Thus $Y\subseteq J$.
Choose a path through $Y$ connecting $C$ to $C'$; its length is at most $t-1$.
Treat each direct edge between $C$ and $C'$ as another connecting path, of length one.
All these paths are internally disjoint.
Their endpoints are distinct, since each cycle vertex has at most one edge outside its cycle.

Put $a=\lfloor D/8\rfloor\ge D/16$.
Partition each cycle into at most $\lceil L/a\rceil\le17L/D$ consecutive arcs, each containing at most $a$ vertices.
Two connecting paths with endpoints in the same pair of arcs would give a cycle in $J$ of length at most
\[
 2t+2(a-1)\le2(t+2a)\le D.
\]
This contradicts the girth of $J$.
There is at most one path per pair of arcs, so the number of these links is at most $289(L/D)^2$.
The stated bound uses the harmless larger coefficient $578$.
Adding the four types of links proves the result.
\end{proof}

\begin{lemma}[The moment at a strict half threshold]
\label{lem:sharp-half-moment}
Use the preceding graph and marked set.
Let $\mathcal C$ be all cycles of $J$ of length at most $L$, and put
\[
 s=\lfloor(D-1)/2\rfloor,\qquad
 a_C=2^{-|C|},\qquad \mu=\sum_{C\in\mathcal C}a_C.
\]
Suppose every disjoint pair has at most $b$ links, and set
\[
 \eta=\frac32\,2^b L^3 2^{-s}.
\]
If $\mu>0$, the total weight $W_2$ of ordered disjoint pairs with exactly two links satisfies
\[
 W_2>\mu^2-(1+\eta)\mu.
\]
\end{lemma}

\begin{proof}
For each cycle use the variable $X_C$ from Section~\ref{sec:polynomial-foundations}, and put $Z=\sum_CX_C$.
Then $\mathbb EZ=\mu$.
The diagonal contribution to its second moment is at most $\mu$.
Intersecting distinct cycles contribute zero.
The pair law and the exceptional-partner estimate give
\[
 \mathbb EZ^2\le\mu^2+W_2+(1+\eta)\mu.
\]
Every cycle in $\mathcal C$ avoids $P$, so the forest event implies $Z=0$.
Its probability exceeds one half.
Consequently Cauchy--Schwarz gives
\[
 \mu^2\le\Pr(Z>0)\mathbb EZ^2<\frac12\mathbb EZ^2.
\]
Substitute the second-moment bound and rearrange.
\end{proof}

\begin{lemma}[A long nested chain]
\label{lem:sharp-nested-chain}
Under the preceding hypotheses, assume
\[
 \mu>20,\qquad \eta\le\frac12,\qquad L^2 2^{-s}\le\frac12.
\]
There are a cycle $C\subseteq J$ of length at most $L$, a neighbor $u$ of $C$, and a family of vertex-disjoint cycles $C_1,\ldots,C_T$ in $J$ such that
\[
 T\ge\frac{8\,2^s}{L^3}.
\]
All $C_i$ lie in the component of $G-V(C)$ containing $P$ and separate $u$ from $P$ there.
They admit an ordering for which their $u$-components $U_i$ in $G-V(C)-V(C_i)$ satisfy
\[
 U_i\subseteq U_{i+1},\qquad V(C_i)\subseteq U_{i+1}.
\]
\end{lemma}

\begin{proof}
For disjoint cycles $C,C'$, call $C'$ far behind $C$ if it lies outside the $P$-component of $G-V(C)$.
This relation is antisymmetric.
To see this, take a path from $P$ to its first vertex on $C$.
The path avoids the far component and hence avoids $C'$.
Connectedness of $C$ then puts all of $C$ in the $P$-component after deleting $C'$.
Thus ordered far pairs have total weight at most $\mu^2/2$.

For fixed $C$, cycles joined to it by an edge meet one of at most $L$ neighboring vertices.
The local weight bound makes their total weight at most $L^2 2^{-s}$.
Ordered adjacent pairs therefore have total weight at most $\mu L^2 2^{-s}$.
After excluding far and adjacent pairs, the preceding moment leaves ordered two-link pairs of total weight greater than
\[
 \frac12\mu^2-2\mu>8\mu.
\]
Weighted averaging gives a cycle $C$ whose remaining partners have total weight greater than eight.

Such a partner $C'$ has no edge to $C$ and has exactly two exterior linking components.
At least one avoids $P$ and contains a neighbor $u$ of $C$.
In $H=G-V(C)$, that partner separates $u$ from $P$.
There are at most $L$ choices of $u$.
For one choice the separating partners have total weight at least $8/L$.

Choose a maximal vertex-disjoint subfamily of these partners, of size $T$.
Its at most $TL$ vertices meet every partner in this family.
The local weight bound gives
\[
 \frac8L\le TL\cdot L2^{-s},
\]
which is the asserted lower bound.

Finally, disjoint connected separators of two connected sets have a consistent order.
Here is a direct justification in the present setting.
For one separator $A$, a path from $u$ to $P$ can be rerouted within $A$ between its first and last visits.
A disjoint connected separator $B$ must therefore lie in the $u$-component or the $P$-component of $H-A$.
If it lay in another component, that rerouted path would avoid $B$.
If $B$ lies on the $P$-side of $A$, then $A$ and the entire $u$-component of $H-A$ lie in the $u$-component of $H-B$.
The symmetric alternative reverses their order.
This orders every pair by inclusion of the $u$-components, giving the claimed nesting.
\end{proof}

% Audited mathematical working proof: root, four_regular, and path_stability passed full independent line audits.
\section{A bootstrap to the unrestricted sharp-order error}
\label{sec:unrestricted-sharp-bootstrap}

This audited mathematical working proof combines the excess-free girth-raising theorem with two conditional clean-block estimates.
The original qualitative marked-cycle estimate suffices through localization and affine-forest bounds.
The independently established polynomial forest theorem supplies a stronger version of that input.
Thus the argument does not assume the sharp conclusion that it seeks to prove.
All logarithms have base two.

\begin{theorem}[An exact half bound above a fixed girth]
\label{thm:bootstrap-exact-half-forest}
There are absolute integers $D_*,R$ such that the following holds.
Let $G$ be a connected cubic multigraph of rank $r\ge R$.
Let $P$ be connected or empty, with $|P|\le3\log r$.
If
\[
 \operatorname{girth}(G-P)\ge D_*,
\]
then
\[
 \Pr(\xi_G-P\text{ is a forest})\le\frac12.
\]
Here an acyclic graph has infinite girth.
\end{theorem}

\begin{proof}
Take $D_*$ from Theorem~\ref{thm:four-excess-free-girth-raising}.
If $P$ is empty, the only even forest is the empty state, of probability $2^{-r}$.
Assume $P$ is nonempty and, for contradiction, that the forest probability exceeds one half.
Write $x=\log r$, $z=\log x$, and put
\[
 A=6e^{168},\qquad
 L=\lfloor Ax\rfloor,\qquad
 t=\lfloor z-2\log z\rfloor.
\]
Use the common forward dyadic sequence $2^jD_*$.
Let $D_0$ be its first value at least $\sqrt x$, and let $D_1$ be its first value at least $x/16$.
For sufficiently large $r$,
\[
 \sqrt x\le D_0<2\sqrt x,\qquad
 x/16\le D_1<x/8,\qquad D_*<D_0<D_1.
\]
Apply the maximal-packing recurrence from the girth-raising theorem along this one sequence.
It supplies nested sets $S_0\subseteq S_1$, disjoint from $P$, such that
\[
 J_j=G-(P\cup S_j)
\]
has girth greater than $D_j$, and
\begin{align*}
 \log(|S_0|+1)
 &\le2D_0+K\log^2(D_0+2)=o(x),\\
 \log(|S_1|+1)
 &\le2D_1+K\log^2(D_1+2)
 \le x/4+o(x).
\end{align*}
In particular,
\[
 |S_0|=r^{o(1)},\qquad |S_1|\le r^{1/4+o(1)}=o(r/x).
\]
Each $S_j$ is the vertex union of a maximal disjoint family of cycles of length at most $D_j$ outside $P$:
the construction only adds disjoint cycles, and its final complement contains no further such cycle.

We use the non-backtracking trace and local cycle-weight estimates from Section~\ref{sec:polynomial-foundations}.
The link and moment estimates are Lemmas~\ref{lem:sharp-uniform-links} and~\ref{lem:sharp-half-moment}.
Their required parameters are verified here.
Write
\[
 s_1=\lfloor(D_1-1)/2\rfloor\ge x/32-2.
\]
The graph $J_1$ is simple, has order $v_1=2r+o(r)$, and has degree deficit
$d_1\le3(|P|+|S_1|)$.
Consequently $Ld_1/(2v_1)=o(1)$ and $L^2 2^{-s_1}=o(1)$.
The even trace sum over $6x\le\ell\le L$ therefore gives
\[
 \mu_1:=
 \sum_{\substack{C\subseteq J_1\\|C|\le L}}2^{-|C|}
 \ge\left(\frac14-o(1)\right)\ln(A/6)>20.
\]

The large-scale link estimate applies because $8\le D_1\le L$ and
$2(t+2\lfloor D_1/8\rfloor)\le D_1$ for large $r$.
The marked set is connected, so it gives the bound
\[
 b_1\le
 1+\frac{2L}{t+1}+2^{t+1}+578(L/D_1)^2
 \le\frac{3Ax}{z}
\]
on the number of links between disjoint cycles of length at most $L$ in $J_1$.
The exceptional-pair error is
\[
 \eta_1\le\frac32\,2^{b_1}L^3 2^{-s_1}
 \le r^{-1/32+o(1)}=o(1).
\]
Only the first refined moment conclusion is needed:
\[
 W_{2,1}>\mu_1^2-(1+\eta_1)\mu_1.
\]
It follows from forest probability greater than one half, without a lower bound on that excess.
No moment estimate is used at scale $D_0$.

Apply Lemma~\ref{lem:sharp-nested-chain}, with the connected set $P$ as anchor.
It gives an auxiliary cycle $C\subseteq J_1$ of length at most $L$, a neighbor $u$ of $C$, and a nested chain
$C_1,\ldots,C_T$ of vertex-disjoint cycles in the component of $G-V(C)$ containing $P$.
Each chain cycle separates $u$ from $P$, and
\[
 T\ge\frac{8\,2^{s_1}}{L^3}
 \ge\frac{2}{A^3}\frac{r^{1/32}}{x^3}.
\]
For reference, the moment inequality leaves partner weight at least eight after discarding ordered far pairs and adjacent pairs.
Choosing one of at most $L$ neighbors of $C$ leaves weight at least $8/L$.
The local cycle-weight bound $L2^{-s_1}$ and maximal vertex-disjointness give the displayed chain length.

Let $H=G-V(C)$, and let $U_i$ be the $u$-component of $H-V(C_i)$.
The nesting gives $U_i\subseteq U_{i+1}$ and $C_i\subseteq U_{i+1}$.
Put $q=\lceil128A^2x\rceil+2$.
Choose disjoint index intervals $[a,b]$ with $b-a=q$, leaving a separator between successive intervals.
For each one, take the full block
\[
 R(a,b)=U_b\setminus\bigl(U_a\cup V(C_a)\bigr).
\]
These blocks are pairwise disjoint.
Each contains all $q-1$ intermediate chain cycles.
Since each of those cycles has length greater than $D_1$,
\[
 |R(a,b)|\ge(q-1)(D_1+1)\ge8L^2
\]
for sufficiently large $r$.
Every edge of $H$ leaving a block ends on $C_a$ or $C_b$, so its boundary in $H$ has at most $2L$ edges.

Discard blocks meeting a neighbor of $C$, a vertex of $S_0$, or a vertex of $P$.
This removes at most $L+|S_0|+1$ blocks.
The last bound uses connectedness of $P$: it cannot cross a block boundary, since that boundary ends on chain cycles disjoint from $P$.
The total number of blocks is of order at least $r^{1/32}/x^4$.
Thus at least two remain after the $r^{o(1)}$ discards.

Every retained block avoids $P\cup S_0$, has girth greater than $D_0$, and has boundary at most $2L$ in $G$.
The boundary statement follows because the block meets no neighbor of the removed auxiliary cycle $C$.
Finally,
\[
 \frac{D_0}{\log(L+2)}\longrightarrow\infty.
\]
Proposition~\ref{prop:path-two-clean-blocks} therefore bounds the original forest probability by $9/25$.
This contradicts the assumption that it exceeds one half and proves the theorem.
\end{proof}

\begin{theorem}[The unrestricted sharp-order upper bound]
\label{thm:bootstrap-unrestricted-sharp-error}
There is an absolute constant $K$ such that every finite graph of cycle rank $r\ge1$ satisfies
\[
 c(G)\le2^{r-1}+Kr2^{r/2}.
\]
The order of the second term is necessary on an infinite sequence of ranks.
\end{theorem}

\begin{proof}
Join distinct components by bridges.
This preserves the cycle count and the rank $r=|E|-|V|+\operatorname{comp}$, so it suffices to prove the connected case.

We use two established inputs.
First, the extremal cycle count $\psi(r)$ is attained by a finite connected cubic multigraph for $r\ge2$.
Attainment starts from a finite graph because the nonempty set of possible cycle counts is a set of integers bounded by $2^r-1$.
The reduction deletes leaves, suppresses degree-two vertices, and replaces each higher-degree vertex by a trivalent tree.
It preserves rank and injects the original cycles.
Multigraphs and simple graphs have the same extremal value because subdivision preserves rank and cycle count.
Second, Theorem~\ref{thm:sharp-order-four-cut-error} gives the claimed order, with an absolute constant $K_4$, for loopless cubic multigraphs with a cyclic cut of size at most four.

Choose an absolute rank threshold $R_0$ above that in
Theorem~\ref{thm:bootstrap-exact-half-forest} and above the fixed-girth deficit thresholds for every girth below $D_*$.
The latter thresholds exist by Corollary~\ref{cor:path-unbounded-girth-reduction}.
Increase $R_0$ if needed for the cubic Moore bound $g\le3\log r$.
Choose $K\ge K_4$ large enough that the asserted bound holds for every rank below $R_0$.
This is possible using the elementary bound $\psi(r)\le2^r-1$.
We prove the theorem by strong induction on the rank.

Let $r\ge R_0$, and choose a cubic multigraph attaining $\psi(r)$.
If it has a loop $e$, that loop is the only cycle containing $e$.
Deleting $e$ leaves a connected graph of rank $r-1$, so the induction hypothesis gives
\[
 c(G)=c(G-e)+1
 \le2^{r-2}+K(r-1)2^{(r-1)/2}+1
 \le2^{r-1}+Kr2^{r/2}.
\]
Assume henceforth that $G$ is loopless.

If it has a cyclic cut of size at most four, apply the established bound with constant $K_4$.
Otherwise its girth is at least five at sufficiently large rank:
a cycle of length at most four has boundary at most four and leaves a cyclic exterior once the rank exceeds an absolute constant.
If its girth is below $D_*$, the fixed-girth deficit gives
$c(G)<2^{r-1}$.
These cases are complete.

In the remaining case, choose a shortest cycle $W$.
It has at most $3\log r$ vertices, and $G-V(W)$ has girth at least $D_*$.
Theorem~\ref{thm:bootstrap-exact-half-forest}, with $P=V(W)$, gives forest probability at most one half.
For any edge $e$ of $W$, each cycle containing $e$ gives two distinct even states by toggling $W$.
Both restrictions outside $V(W)$ are forests.
Consequently the number $c_e$ of cycles containing $e$ satisfies
\[
 2c_e\le\frac12\,2^r,
 \qquad c_e\le2^{r-2}.
\]
The graph $G-e$ is connected and has rank $r-1$.
Apply the induction hypothesis to this general connected graph:
\[
 c(G)=c(G-e)+c_e
 \le2^{r-1}+K(r-1)2^{(r-1)/2}
 \le2^{r-1}+Kr2^{r/2}.
\]
This completes the induction.
It does not require $G-e$ to remain cubic.

Finally, the cyclic-square family from the sharp four-cut theorem has ranks $r=2q+1$ and excess
\[
 \left(2q+\frac12\right)2^q-3q
 =\Theta(r2^{r/2}).
\]
Thus the error term cannot be replaced uniformly by $o(r2^{r/2})$.
\end{proof}


\section{The matching lower-order examples}
For completeness, we give the exact count in the known cyclic-square family~\cite{HK}.
For $q\ge2$, use indices modulo $q$.
Block $i$ has vertices $a_i,b_i,c_i,d_i$ and all four edges between $\{a_i,b_i\}$ and $\{c_i,d_i\}$.
Join consecutive blocks by the edges $c_i a_{i+1}$ and $d_i b_{i+1}$.
The resulting simple cubic graph $G_q$ has $4q$ vertices, $6q$ edges, and rank $2q+1$.

A cycle uses either one edge at every interface, or zero or two at every interface.
In the first case there are $2^q$ choices of interfaces and two local paths in every square, giving $4^q$ cycles.
If no interface is used, there are $q$ square cycles.
If a nonempty proper subset of interfaces is used, connectedness makes it a cyclic interval.
For an interval of $d$ interfaces, the two end blocks each have two caps and each inner block has two matchings.
There are $q2^{d+1}$ such cycles for each $1\le d\le q-1$.

When all interfaces are used twice, each square supplies one of two matchings.
Their composition must be the transposition to give one cycle.
This gives $2^{q-1}$ cycles.
Adding the cases yields
\[
 \begin{aligned}
 c(G_q)
 &=4^q+q+q\sum_{d=1}^{q-1}2^{d+1}+2^{q-1}\\
 &=2^{2q}+\left(2q+\frac12\right)2^q-3q.
 \end{aligned}
\]
With $r=2q+1$, the excess is $\Theta(r2^{r/2})$.
Therefore no uniform $o(r2^{r/2})$ error can replace the proved upper order.

\section{Verification and formalization status}
The proof does not infer an asymptotic theorem from finite enumeration.
Exact checks instead test the finite identities at the new interfaces between the arguments.
The iterative-packing checker verifies 5,572 selected-star sets, 636 middle intervals, and 114 disjoint segment pairs.
It also checks 4,482 core prunings and all 4,163 conditional affine fibers from 421 disjoint-core pairs.
The fibers are checked as complete Cartesian products, including adjacent cores.

The localization checker compares the complete restriction laws in 1,193 cases, including 265 disconnected regions and 56 padded cap cycles.
It enumerates 93,632 even states.
The shortest-cycle checker verifies 292 prefix deletions and the associated path partitions.
The cyclic-square formula has an independent cycle-space check for $2\le q\le7$.
The earlier spectral and parallel-link checks are preserved with their source packets.

A separate verification record lists the final source hashes, build results, check records, and audit files:
\begin{center}
\nolinkurl{verification-unrestricted-sharp-2026-11-11.json}.
\end{center}
All new work remains separate from the original manuscript and the root Lean imports.
Fable's existing Lean modules cover several earlier link, moment, separation, and transfer steps.
They do not yet formalize the new affine bounds, iterative packing, or complete sharp-order assembly.

The shared \nolinkurl{CODEX-TO-FABLE-LIVE.txt} records the partial results and the completed mathematical argument.
The new proof settles the requested error order within this research programme.
It does not determine the optimal absolute constant or the exact value of $\psi(r)$.

\begin{thebibliography}{9}
\bibitem{HK}
P. Horak and D. Kim,
\emph{Connected cubic graphs with the maximum number of perfect matchings},
Journal of Graph Theory 99 (2022), no. 4, 671--690.
Updated author version: \url{https://arxiv.org/abs/2006.13459}.
\end{thebibliography}
\end{document}
